% Content of Mock exam 3, shared by the problem sheet and the solution sheet.
% Solutions are wrapped in the `solution` environment and appear only when the
% driver defines \ipsshowsolutions.

% Marks per part, printed in the accent colour.
\newcommand{\pts}[1]{{\normalfont\small\bfseries\color{ipsAccent}[#1~pts]}\enspace}

\thispagestyle{fancy}

\begingroup
\noindent{\color{ipsAccent}\rule{\linewidth}{1.2pt}}\\[8pt]
{\LARGE\bfseries\color{ipsPrimary} Mock exam 3\ifipssolutions: full solutions\fi\par}
\vskip7pt
{\normalsize\color{ipsAccent} Planetary Systems final exam practice \textbullet\ closed book \textbullet\ 120 minutes\par}
\vskip3pt
{\footnotesize Planetary Systems (WBAS002-05) \textbullet\ Kapteyn Astronomical Institute, University of Groningen\par}
\vskip7pt
\noindent{\color{ipsAccent}\rule{\linewidth}{0.6pt}}
\vskip10pt
\endgroup

This mock exam has the same shape and level as the final exam: 6 questions of 15 points each, 90 points in total, in a 120-minute closed-book sitting where a calculator and a pen are all that is allowed.
Your grade is $1 + \text{points}/10$; the pass mark of 6.0 corresponds to 50 points.
The twelve relations on the exam formula list are assumed known; every other formula a question needs is printed in the question itself.
Marks per part are shown in brackets, and the steps carry marks, not only the number: show your algebra.
Some parts state an intermediate result (``show that\,\ldots''); use it to check your work, and to continue if a step resists.
Carry at least four significant figures through the intermediate steps and round only at the end.

\begin{given}[Constants and data]
$G = 6.674 \times 10^{-11}\ \mathrm{m^3\,kg^{-1}\,s^{-2}}$ \quad
$1\ \mathrm{AU} = 1.496 \times 10^{11}\ \mathrm{m}$ \quad
$1\ \mathrm{yr} = 3.156 \times 10^{7}\ \mathrm{s} = 365.25\ \mathrm{d}$ \quad
$1\ \mathrm{d} = 86\,400\ \mathrm{s}$ \quad
$\kB = 1.381 \times 10^{-23}\ \mathrm{J\,K^{-1}}$ \quad
$m_u = 1.661 \times 10^{-27}\ \mathrm{kg}$ \quad
$\sigma = 5.670 \times 10^{-8}\ \mathrm{W\,m^{-2}\,K^{-4}}$ \quad
$S_0 = 1361\ \mathrm{W\,m^{-2}}$ (solar constant at $1\ \mathrm{AU}$)

\vskip6pt
\begin{tabular}{@{}llp{7.2cm}@{}}
\toprule
Body & $a$, $A$ & Further data \\
\midrule
Sun & & $R = 6.957 \times 10^{5}\ \mathrm{km}$ \\
Jupiter & & $R = 69\,911\ \mathrm{km}$ \\
Callisto & & orbits Jupiter with $a = 1.883 \times 10^{6}\ \mathrm{km}$, $P = 16.69\ \mathrm{d}$ \\
Earth & $1.000\ \mathrm{AU}$, $A = 0.30$ & $P = 1\ \mathrm{yr}$, $R = 6371\ \mathrm{km}$, $g = 9.81\ \mathrm{m\,s^{-2}}$, surface $T = 288\ \mathrm{K}$, mean molecular weight $\mu = 28.97$ \\
Moon & & $g = 1.62\ \mathrm{m\,s^{-2}}$ \\
Mercury & $0.387\ \mathrm{AU}$, $A = 0.088$ & $R = 2440\ \mathrm{km}$, $\bar{\rho} = 5427\ \mathrm{kg\,m^{-3}}$ \\
Titan & & surface $T = 94\ \mathrm{K}$, surface pressure $1.5 \times 10^{5}\ \mathrm{Pa}$, methane mole fraction $5\%$ \\
\bottomrule
\end{tabular}

\vskip4pt
{\footnotesize A total power $L$ spread over a sphere of radius $r$ gives a flux $L/(4\pi r^2)$; for sunlight this means $S = S_0/a^2$ with $a$ in AU.
A body of mass $m$ moving at speed $v$ carries the kinetic energy $E = \tfrac{1}{2} m v^2$.
Dry air on Earth cools by $9.8\ \mathrm{K}$ per kilometre of ascent (the dry adiabatic lapse rate).}
\end{given}


% ════════════════════════════════════════════════════════════
\problem{Weighing the Sun and Jupiter}

An orbit is a balance you can read a mass from: the central mass alone sets how fast a small companion circles at a given distance.
This problem uses only relations from the formula list.

\parti \pts{5} Use the Earth's orbit (data table) to compute the mass of the Sun in kilograms.

\begin{solution}
Kepler's third law solved for the central mass, with $a = 1.496 \times 10^{11}\ \mathrm{m}$ and $P = 3.156 \times 10^{7}\ \mathrm{s}$:
\[
  M_\odot = \frac{4\pi^2 a^3}{G P^2}
          = \frac{4\pi^2 \times (1.496 \times 10^{11})^3}{6.674 \times 10^{-11} \times (3.156 \times 10^{7})^2}
          = \frac{1.3218 \times 10^{35}}{6.6475 \times 10^{4}}
          = 1.988 \times 10^{30}\ \mathrm{kg}.
\]
\answer{$M_\odot = 1.988 \times 10^{30}\ \mathrm{kg}$, the accepted value to the precision of the input data.}
\end{solution}

\parti \pts{4} Use Callisto's orbit (data table) to compute the mass of Jupiter in kilograms.

\begin{solution}
The same relation with $a = 1.883 \times 10^{9}\ \mathrm{m}$ and $P = 16.69 \times 86\,400 = 1.442 \times 10^{6}\ \mathrm{s}$:
\[
  M_{\mathrm{J}} = \frac{4\pi^2 \times (1.883 \times 10^{9})^3}{6.674 \times 10^{-11} \times (1.442 \times 10^{6})^2}
                 = \frac{2.6358 \times 10^{29}}{1.3878 \times 10^{2}}
                 = 1.899 \times 10^{27}\ \mathrm{kg}.
\]
\answer{$M_{\mathrm{J}} = 1.899 \times 10^{27}\ \mathrm{kg}$, about $10^{-3}$ of the solar mass.}
\end{solution}

\parti \pts{3} Compute the mean density of the Sun and of Jupiter.

\begin{solution}
From $\bar{\rho} = M / (\tfrac{4}{3}\pi R^3)$ with $R_\odot = 6.957 \times 10^{8}\ \mathrm{m}$ and $R_{\mathrm{J}} = 6.9911 \times 10^{7}\ \mathrm{m}$:
\[
  \bar{\rho}_\odot = \frac{1.988 \times 10^{30}}{\tfrac{4}{3}\pi (6.957 \times 10^{8})^3} = 1409\ \mathrm{kg\,m^{-3}},
  \qquad
  \bar{\rho}_{\mathrm{J}} = \frac{1.899 \times 10^{27}}{\tfrac{4}{3}\pi (6.9911 \times 10^{7})^3} = 1327\ \mathrm{kg\,m^{-3}}.
\]
\answer{$\bar{\rho}_\odot = 1409\ \mathrm{kg\,m^{-3}}$ and $\bar{\rho}_{\mathrm{J}} = 1327\ \mathrm{kg\,m^{-3}}$: both close to the density of water.}
\end{solution}

\parti \pts{3} Judge this claim in 2 to 3 sentences, using your result from part (c):
\emph{``Because their mean densities are nearly equal, Jupiter and the Sun must have very similar interiors.''}

\begin{solution}
False.
The similar mean densities reflect a similar hydrogen-helium composition, but the interiors differ fundamentally: the Sun is a plasma at millions of kelvin supported by thermal pressure from fusion, while Jupiter is a cold body of molecular and metallic hydrogen supported by the pressure of compressed, partially degenerate matter.
A mean density is a single number that erases the internal structure; bodies with very different temperature and pressure profiles can share it.
\answer{False: the near-equal mean densities come from similar composition, not similar structure; a fusing plasma ball and a cold compressed fluid planet merely happen to average out alike.}
\end{solution}


% ════════════════════════════════════════════════════════════
\problem{Reading a cratered surface}

Impact craters are both a hazard record and a clock.
In the gravity regime, an impact of energy $E$ into a target of density $\rho_t$ and surface gravity $g$ opens a crater of diameter
\[
  D \approx \left( \frac{E}{\rho_t\, g} \right)^{1/4}
\]
(all quantities in SI units, $D$ in metres).
Craters larger than a transition diameter $D_{\mathrm{t}}$ collapse into complex craters with central peaks and terraces; on the Moon $D_{\mathrm{t}} \approx 15\ \mathrm{km}$, and $D_{\mathrm{t}}$ scales inversely with surface gravity, $D_{\mathrm{t}} \propto 1/g$.
For surfaces younger than about $3\ \mathrm{Gyr}$, the cumulative density of craters larger than $1\ \mathrm{km}$ grows linearly with age $T$:
\[
  N(\geq 1\ \mathrm{km}) \approx b\,T, \qquad b = 8.38 \times 10^{-4}\ \mathrm{km^{-2}\,Gyr^{-1}}.
\]

\parti \pts{4} A rocky asteroid of diameter $1.2\ \mathrm{km}$ and density $3000\ \mathrm{kg\,m^{-3}}$ hits at $18\ \mathrm{km\,s^{-1}}$.
Show that its kinetic energy is $4.397 \times 10^{20}\ \mathrm{J}$.

\begin{solution}
Mass of the impactor (radius $600\ \mathrm{m}$):
\[
  m = \frac{4}{3}\pi r^3 \rho = \frac{4}{3}\pi \times (600)^3 \times 3000 = 2.714 \times 10^{12}\ \mathrm{kg}.
\]
Kinetic energy at $v = 1.8 \times 10^{4}\ \mathrm{m\,s^{-1}}$:
\[
  E = \frac{1}{2} m v^2 = \frac{1}{2} \times 2.714 \times 10^{12} \times (1.8 \times 10^{4})^2 = 4.397 \times 10^{20}\ \mathrm{J}.
\]
\answer{$E = 4.397 \times 10^{20}\ \mathrm{J}$, roughly two thousand times the largest nuclear test.}
\end{solution}

\parti \pts{4} This asteroid strikes the Moon, where the target rock has density $2500\ \mathrm{kg\,m^{-3}}$.
Compute the crater diameter and state whether the crater is simple or complex.

\begin{solution}
With $E = 4.397 \times 10^{20}\ \mathrm{J}$, $\rho_t = 2500\ \mathrm{kg\,m^{-3}}$, $g = 1.62\ \mathrm{m\,s^{-2}}$:
\[
  D \approx \left( \frac{4.397 \times 10^{20}}{2500 \times 1.62} \right)^{1/4}
    = \left( 1.0857 \times 10^{17} \right)^{1/4}
    = 1.815 \times 10^{4}\ \mathrm{m} = 18.15\ \mathrm{km}.
\]
This exceeds the lunar transition diameter of $15\ \mathrm{km}$.
\answer{$D \approx 18.2\ \mathrm{km}$, larger than $D_{\mathrm{t}} = 15\ \mathrm{km}$: a complex crater with a central peak and terraced walls.}
\end{solution}

\parti \pts{3} The same asteroid strikes the Earth, where the target rock has density $2700\ \mathrm{kg\,m^{-3}}$.
Compute the crater diameter and the transition diameter on Earth, and classify the crater.

\begin{solution}
On Earth, $g = 9.81\ \mathrm{m\,s^{-2}}$:
\[
  D \approx \left( \frac{4.397 \times 10^{20}}{2700 \times 9.81} \right)^{1/4}
    = \left( 1.6601 \times 10^{16} \right)^{1/4}
    = 1.135 \times 10^{4}\ \mathrm{m} = 11.35\ \mathrm{km}.
\]
The transition diameter scales down with gravity:
\[
  D_{\mathrm{t,Earth}} = 15 \times \frac{1.62}{9.81} = 2.477\ \mathrm{km}.
\]
\answer{$D \approx 11.4\ \mathrm{km}$ against $D_{\mathrm{t}} \approx 2.5\ \mathrm{km}$: strongly complex. The same impactor makes a smaller crater on Earth because the stronger gravity resists excavation, yet the crater collapses more readily.}
\end{solution}

\parti \pts{4} A lava plain on the Moon shows $N(\geq 1\ \mathrm{km}) = 1.8 \times 10^{-3}\ \mathrm{km^{-2}}$.
Compute its age, check that the linear relation applies, and explain in 1 to 2 sentences why crater counting cannot date the oldest lunar highlands this way.

\begin{solution}
\[
  T = \frac{N}{b} = \frac{1.8 \times 10^{-3}}{8.38 \times 10^{-4}} = 2.148\ \mathrm{Gyr}.
\]
This is below $3\ \mathrm{Gyr}$, so the linear (constant-flux) segment of the crater production rate applies.
The oldest highlands date from before about $3.6\ \mathrm{Gyr}$, when the impact flux was far higher and rose steeply back in time: there the crater density grows exponentially with age rather than linearly, and surfaces begin to saturate, so the simple proportionality fails.
\answer{$T \approx 2.1\ \mathrm{Gyr}$; for the ancient highlands the early bombardment breaks the linear flux assumption, so the relation printed above cannot be extrapolated there.}
\end{solution}


% ════════════════════════════════════════════════════════════
\problem{The birth of a core and a dynamo}

Terrestrial planets separate into an iron core and a rocky mantle early in their history, and a convecting liquid core can drive a magnetic dynamo.
In a magma ocean of viscosity $\mu$, an iron droplet of radius $r$ and excess density $\Delta\rho$ sinks at the Stokes settling speed
\[
  v = \frac{2\, \Delta\rho\, g\, r^2}{9 \mu}.
\]
A convecting core of flow speed $U$, size $L$, and magnetic diffusivity $\eta$ sustains a dynamo when the magnetic Reynolds number
\[
  \mathrm{Rm} = \frac{U L}{\eta}
\]
exceeds a critical value of order $10$ to $100$.

\parti \pts{4} The radioactive isotope $^{182}\mathrm{Hf}$ decays to $^{182}\mathrm{W}$ with a half-life of $8.9\ \mathrm{Myr}$ and is the main clock for dating core formation.
Compute the decay constant $\lambda$, the fraction of the initial $^{182}\mathrm{Hf}$ remaining $30\ \mathrm{Myr}$ after the birth of the solar system, and explain in 1 to 2 sentences why this clock cannot date events later than about $45\ \mathrm{Myr}$.

\begin{solution}
\[
  \lambda = \frac{\ln 2}{t_{1/2}} = \frac{0.6931}{8.9} = 7.788 \times 10^{-2}\ \mathrm{Myr^{-1}}.
\]
At $t = 30\ \mathrm{Myr}$:
\[
  \frac{N}{N_0} = e^{-\lambda t} = e^{-7.788 \times 10^{-2} \times 30} = e^{-2.3364} = 9.668 \times 10^{-2} \approx 9.7\%.
\]
By $45\ \mathrm{Myr}$, about five half-lives have passed and only about $3\%$ of the parent remains; after that, essentially all $^{182}\mathrm{Hf}$ has decayed, so later core-mantle separation no longer changes the tungsten isotope record.
\answer{$\lambda = 7.788 \times 10^{-2}\ \mathrm{Myr^{-1}}$; $9.7\%$ remains at $30\ \mathrm{Myr}$. A radioactive clock goes silent after roughly five half-lives, here $\approx 45\ \mathrm{Myr}$.}
\end{solution}

\parti \pts{4} In a crystal-rich magma ocean with $g = 5\ \mathrm{m\,s^{-2}}$ and effective viscosity $\mu = 100\ \mathrm{Pa\,s}$, an iron droplet of radius $1\ \mathrm{cm}$ has excess density $\Delta\rho = 4000\ \mathrm{kg\,m^{-3}}$.
Compute its settling speed and the time it needs to sink through a magma ocean $800\ \mathrm{km}$ deep.

\begin{solution}
\[
  v = \frac{2 \times 4000 \times 5 \times (0.01)^2}{9 \times 100} = 4.444 \times 10^{-3}\ \mathrm{m\,s^{-1}}.
\]
Sinking time over $8 \times 10^{5}\ \mathrm{m}$:
\[
  t = \frac{8 \times 10^{5}}{4.444 \times 10^{-3}} = 1.800 \times 10^{8}\ \mathrm{s} \approx 5.7\ \mathrm{yr}.
\]
\answer{$v = 4.444 \times 10^{-3}\ \mathrm{m\,s^{-1}}$; the droplet crosses the magma ocean in about $6$ years. Once a planet is molten, metal-silicate separation is essentially instantaneous on geological timescales, which is why cores form within the first tens of Myr.}
\end{solution}

\parti \pts{4} Early Mars had a convecting liquid core with $U = 5 \times 10^{-4}\ \mathrm{m\,s^{-1}}$, $L = 1.8 \times 10^{6}\ \mathrm{m}$, and $\eta = 1\ \mathrm{m^2\,s^{-1}}$.
Compute the magnetic Reynolds number and state whether a dynamo is expected then, and what happens once core convection stops.

\begin{solution}
\[
  \mathrm{Rm} = \frac{5 \times 10^{-4} \times 1.8 \times 10^{6}}{1} = 900,
\]
far above the critical value of $10$ to $100$: a dynamo is expected.
When core convection stops, $U \to 0$ and $\mathrm{Rm}$ falls below critical; the field is no longer regenerated and decays away resistively.
\answer{$\mathrm{Rm} = 900 \gg \mathrm{Rm_c}$: early Mars could sustain a dynamo. A dynamo needs stirring, not just liquid iron; once convection dies, so does the global field.}
\end{solution}

\parti \pts{3} Judge this claim in 2 to 3 sentences:
\emph{``Mars has no global magnetic field today, so it can never have had one.''}

\begin{solution}
False.
Portions of the ancient southern-highlands crust carry strong remanent magnetization, which the rocks could only acquire by cooling in a global field: direct evidence for an early Martian dynamo.
The dynamo died when the small core stopped convecting, but the magnetized crust preserves the memory of the field, just as your part (c) reasoning predicts.
\answer{False: crustal remanent magnetization records an ancient dynamo; the absence of a field today only shows that core convection has ceased.}
\end{solution}


% ════════════════════════════════════════════════════════════
\problem{Cloud bases on Earth and Titan}

Clouds form where a gas becomes saturated.
The saturation vapour pressure over a condensed phase follows the integrated Clausius-Clapeyron relation
\[
  P_{\mathrm{sat}}(T) = P_{\mathrm{ref}} \exp\!\left[ -\frac{L_v}{R_v} \left( \frac{1}{T} - \frac{1}{T_{\mathrm{ref}}} \right) \right],
\]
where $L_v/R_v$ is a constant of the condensing species.
For water, $L_v/R_v = 5400\ \mathrm{K}$ with reference point $(T_{\mathrm{ref}}, P_{\mathrm{ref}}) = (273\ \mathrm{K},\ 611\ \mathrm{Pa})$.
For methane, $L_v/R_v = 980\ \mathrm{K}$ with reference point $(90.7\ \mathrm{K},\ 1.17 \times 10^{4}\ \mathrm{Pa})$.

\parti \pts{3} Compute the pressure scale height of the Earth's atmosphere near the surface (data table).

\begin{solution}
\[
  H = \frac{\kB T}{\mu\, m_u\, g}
    = \frac{1.381 \times 10^{-23} \times 288}{28.97 \times 1.661 \times 10^{-27} \times 9.81}
    = 8426\ \mathrm{m} \approx 8.4\ \mathrm{km}.
\]
\answer{$H \approx 8.4\ \mathrm{km}$: the pressure drops by a factor $e$ for every $8.4\ \mathrm{km}$ of altitude.}
\end{solution}

\parti \pts{4} Compute the saturation vapour pressure of water at $303\ \mathrm{K}$ (a hot summer day).

\begin{solution}
\[
  P_{\mathrm{sat}}(303) = 611 \exp\!\left[ -5400 \left( \frac{1}{303} - \frac{1}{273} \right) \right]
                        = 611\, e^{1.9584}
                        = 611 \times 7.088 = 4331\ \mathrm{Pa}.
\]
\answer{$P_{\mathrm{sat}} \approx 4.33\ \mathrm{kPa}$, seven times the value at the freezing point: the exponential is steep for water because $L_v/R_v$ is large.}
\end{solution}

\parti \pts{4} An air parcel at the surface has temperature $303\ \mathrm{K}$ and water vapour pressure $1.7\ \mathrm{kPa}$.
Compute the dew point of the parcel, and use the dry adiabatic lapse rate to estimate the altitude of the cloud base it forms when it rises.

\begin{solution}
The dew point is where the actual vapour pressure equals saturation, so invert the Clausius-Clapeyron relation for $P_{\mathrm{sat}}(T_{\mathrm{d}}) = 1700\ \mathrm{Pa}$:
\[
  \frac{1}{T_{\mathrm{d}}} = \frac{1}{273} - \frac{1}{5400} \ln\!\frac{1700}{611}
                           = 3.6630 \times 10^{-3} - \frac{1.0233}{5400}
                           = 3.4735 \times 10^{-3}\ \mathrm{K^{-1}},
\]
so $T_{\mathrm{d}} = 287.9\ \mathrm{K}$.
The rising parcel cools at $9.8\ \mathrm{K\,km^{-1}}$ and reaches its dew point after
\[
  z \approx \frac{303 - 287.9}{9.8} = 1.541\ \mathrm{km}.
\]
\answer{$T_{\mathrm{d}} = 287.9\ \mathrm{K}$ and a cloud base near $1.5\ \mathrm{km}$: the flat bottoms of cumulus clouds trace the altitude where rising parcels saturate.}
\end{solution}

\parti \pts{4} On Titan, methane plays the role of water.
Compute the saturation vapour pressure of methane at Titan's surface temperature, compare it with the actual methane partial pressure at the surface (data table), and explain in 1 to 2 sentences what this means for Titan's lakes and clouds.

\begin{solution}
At $T = 94\ \mathrm{K}$:
\[
  P_{\mathrm{sat}}(94) = 1.17 \times 10^{4} \exp\!\left[ -980 \left( \frac{1}{94} - \frac{1}{90.7} \right) \right]
                       = 1.17 \times 10^{4}\, e^{0.3793}
                       = 1.710 \times 10^{4}\ \mathrm{Pa}.
\]
The actual methane partial pressure is $0.05 \times 1.5 \times 10^{5} = 7.5 \times 10^{3}\ \mathrm{Pa}$, so the relative humidity is
\[
  \frac{7.5 \times 10^{3}}{1.710 \times 10^{4}} = 0.44.
\]
The surface air is undersaturated: the methane lakes slowly evaporate into it, and clouds form only where air rises and cools.
Because $L_v/R_v$ is small for methane, the saturation curve is shallow, and methane clouds can form over a wide range of altitudes rather than at one sharp cloud base as for water on Earth.
\answer{$P_{\mathrm{sat}} = 17.1\ \mathrm{kPa}$ against a partial pressure of $7.5\ \mathrm{kPa}$: relative humidity $\approx 44\%$. Titan's lakes evaporate into undersaturated air, and its shallow methane saturation curve spreads clouds over a broad altitude range.}
\end{solution}


% ════════════════════════════════════════════════════════════
\problem{Mercury, the shrinking planet}

Mercury is scarred by long cliffs (lobate scarps) that record a global contraction of the planet as its interior cooled.
A surface element facing the Sun directly reaches the subsolar temperature given by
\[
  \sigma T_{\mathrm{ss}}^4 = S\,(1 - A).
\]
Cooling the whole planet by $\Delta T$ releases the heat $E = M c\, \Delta T$, with specific heat capacity $c = 800\ \mathrm{J\,kg^{-1}\,K^{-1}}$, and shrinks the radius by
\[
  \Delta R = \frac{\alpha_V}{3}\, R\, \Delta T, \qquad \alpha_V = 3 \times 10^{-5}\ \mathrm{K^{-1}}.
\]

\parti \pts{4} Compute the solar flux at Mercury's orbit, its global equilibrium temperature, and its subsolar temperature.

\begin{solution}
\[
  S = \frac{S_0}{a^2} = \frac{1361}{0.387^2} = 9087\ \mathrm{W\,m^{-2}},
\]
about $6.7$ times the flux at Earth.
The global equilibrium temperature with $A = 0.088$:
\[
  T_{\mathrm{eq}} = \left[ \frac{S (1-A)}{4\sigma} \right]^{1/4}
                  = \left[ \frac{9087 \times 0.912}{4 \times 5.670 \times 10^{-8}} \right]^{1/4}
                  = \left( 3.6539 \times 10^{10} \right)^{1/4} = 437.2\ \mathrm{K}.
\]
The subsolar point absorbs the full flux on one face instead of spreading it over the sphere:
\[
  T_{\mathrm{ss}} = \left[ \frac{9087 \times 0.912}{5.670 \times 10^{-8}} \right]^{1/4}
                  = \left( 1.4616 \times 10^{11} \right)^{1/4} = 618.3\ \mathrm{K},
\]
a factor $4^{1/4} = 1.414$ above the global value.
\answer{$S = 9087\ \mathrm{W\,m^{-2}}$, $T_{\mathrm{eq}} = 437\ \mathrm{K}$, $T_{\mathrm{ss}} = 618\ \mathrm{K}$: hot enough at noon to melt lead, yet Mercury's poles hold water ice in permanent shadow.}
\end{solution}

\parti \pts{4} Compute Mercury's mass from its radius and mean density, and show that cooling the whole planet by $\Delta T = 200\ \mathrm{K}$ releases $5.283 \times 10^{28}\ \mathrm{J}$.

\begin{solution}
\[
  M = \frac{4}{3}\pi R^3 \bar{\rho} = \frac{4}{3}\pi \times (2.440 \times 10^{6})^3 \times 5427 = 3.302 \times 10^{23}\ \mathrm{kg}.
\]
The released heat:
\[
  E = M c\, \Delta T = 3.302 \times 10^{23} \times 800 \times 200 = 5.283 \times 10^{28}\ \mathrm{J}.
\]
\answer{$M = 3.302 \times 10^{23}\ \mathrm{kg}$ and $E = 5.283 \times 10^{28}\ \mathrm{J}$.}
\end{solution}

\parti \pts{4} Assume this heat left the planet at a constant rate over $4.5\ \mathrm{Gyr}$.
Compute the mean surface heat flow in $\mathrm{mW\,m^{-2}}$ and compare it with the Earth's present mean surface heat flow of $87\ \mathrm{mW\,m^{-2}}$.

\begin{solution}
Over $4.5\ \mathrm{Gyr} = 4.5 \times 10^{9} \times 3.156 \times 10^{7} = 1.4202 \times 10^{17}\ \mathrm{s}$:
\[
  P = \frac{5.283 \times 10^{28}}{1.4202 \times 10^{17}} = 3.720 \times 10^{11}\ \mathrm{W}.
\]
Spread over the surface $4\pi R^2 = 7.4815 \times 10^{13}\ \mathrm{m^2}$:
\[
  q = \frac{3.720 \times 10^{11}}{7.4815 \times 10^{13}} = 4.972 \times 10^{-3}\ \mathrm{W\,m^{-2}} \approx 5.0\ \mathrm{mW\,m^{-2}},
\]
about $17$ times below the Earth's $87\ \mathrm{mW\,m^{-2}}$.
Mercury is small: its large surface-to-volume ratio let it lose its internal heat early, and today only a trickle remains, while the larger Earth still leaks the heat of its formation and of radioactive decay.
\answer{$q \approx 5\ \mathrm{mW\,m^{-2}}$, far below Earth's $87\ \mathrm{mW\,m^{-2}}$: small planets cool fast and finish their geological activity early.}
\end{solution}

\parti \pts{3} Compute the radius change from the $200\ \mathrm{K}$ cooling and compare it with the $7\ \mathrm{km}$ of contraction inferred from Mercury's lobate scarps.

\begin{solution}
\[
  \Delta R = \frac{\alpha_V}{3} R\, \Delta T = \frac{3 \times 10^{-5}}{3} \times 2.440 \times 10^{6} \times 200 = 4880\ \mathrm{m} \approx 4.9\ \mathrm{km}.
\]
This matches the scarp record to within a factor of order unity: secular cooling of a few hundred kelvin explains the observed global contraction.
\answer{$\Delta R \approx 4.9\ \mathrm{km}$, the same order as the $\leq 7\ \mathrm{km}$ read from the scarps: Mercury's wrinkled face is the direct signature of a planet that cooled and shrank.}
\end{solution}


% ════════════════════════════════════════════════════════════
\problem{A planet around a K dwarf}

Stars less massive than the Sun burn slower and live longer.
On the main sequence, luminosity and lifetime scale with stellar mass as
\[
  \frac{L}{L_\odot} = \left( \frac{M}{\Msun} \right)^{3.5},
  \qquad
  t_{\mathrm{MS}} \approx 10 \left( \frac{M}{\Msun} \right)^{-2.5}\ \mathrm{Gyr},
\]
and Kepler's third law for an orbit around a star of mass $M$ reads, in convenient units,
\[
  P^2 = \frac{a^3}{M/\Msun} \qquad (P\ \mathrm{in\ yr},\ a\ \mathrm{in\ AU}).
\]
Consider a K dwarf of mass $M = 0.7\ \Msun$.

\parti \pts{3} Compute the luminosity of the K dwarf in solar units.

\begin{solution}
\[
  \frac{L}{L_\odot} = 0.7^{3.5} = 0.2870.
\]
\answer{$L = 0.287\, L_\odot$: a $30\%$ reduction in mass cuts the luminosity by a factor of $3.5$, the steep price of the mass-luminosity relation.}
\end{solution}

\parti \pts{4} Compute the main-sequence lifetime of the K dwarf, compare it with the Sun's, and state in one sentence what this means for life on its planets.

\begin{solution}
\[
  t_{\mathrm{MS}} = 10 \times 0.7^{-2.5} = 24.39\ \mathrm{Gyr},
\]
against $10\ \mathrm{Gyr}$ for the Sun: almost $2.5$ times longer, and longer than the current age of the universe.
A planet around a K dwarf has a stable stellar energy source for tens of gigayears, giving biology far more time than Earth will get before the Sun leaves the main sequence.
\answer{$t_{\mathrm{MS}} \approx 24\ \mathrm{Gyr}$: no K dwarf that ever formed has yet left the main sequence.}
\end{solution}

\parti \pts{4} Show that a planet receives the same flux as the Earth at $d = 0.5357\ \mathrm{AU}$ from this star, and verify that with the Earth's albedo it has the Earth's equilibrium temperature.

\begin{solution}
The flux matches the Earth's when $L/(4\pi d^2) = L_\odot/(4\pi\, \mathrm{AU}^2)$, so
\[
  d = \sqrt{\frac{L}{L_\odot}}\ \mathrm{AU} = \sqrt{0.2870} = 0.5357\ \mathrm{AU}.
\]
There the stellar flux is $S = S_0 = 1361\ \mathrm{W\,m^{-2}}$ by construction, and with $A = 0.30$:
\[
  T_{\mathrm{eq}} = \left[ \frac{1361 \times 0.70}{4 \times 5.670 \times 10^{-8}} \right]^{1/4}
                  = \left( 4.2006 \times 10^{9} \right)^{1/4} = 254.6\ \mathrm{K},
\]
identical to the Earth's equilibrium temperature, as it must be for identical flux and albedo.
\answer{$d = 0.5357\ \mathrm{AU}$ and $T_{\mathrm{eq}} = 255\ \mathrm{K}$: the habitable zone moves inward around dimmer stars, tracking constant flux.}
\end{solution}

\parti \pts{4} Compute the orbital period of this planet.
Then judge this claim in 2 to 3 sentences:
\emph{``The smaller the star, the better for life, because the star lives longer.''}

\begin{solution}
\[
  P = \sqrt{\frac{0.5357^3}{0.7}} = \sqrt{\frac{0.1537}{0.7}} = \sqrt{0.2196} = 0.4686\ \mathrm{yr} = 171\ \mathrm{d}.
\]
The claim oversimplifies.
Longer lifetimes do favour life, but for the smallest stars (M dwarfs) the habitable zone lies so close in that planets become tidally locked, and young M dwarfs bombard their planets with flares and high-energy radiation that can strip atmospheres.
K dwarfs may offer the best compromise: lifetimes of tens of gigayears with a habitable zone still distant enough to soften those hazards.
\answer{$P = 0.469\ \mathrm{yr} \approx 171\ \mathrm{d}$. Partly true at best: stellar longevity helps, but proximity-driven tidal locking and stellar activity work against the smallest stars.}
\end{solution}
